\documentclass[12pt]{article}
\usepackage[utf8]{inputenc}
\usepackage[T1]{fontenc}
\usepackage{hyperref}
\title{Pokang\'a Language Resources for K-12 Curriculum Development}
\author{Compiled by Claude Sonnet 5.5 (claude-sonnet-5-5) for Andr\'as Kornai}
\date{}
\begin{document}
\maketitle

\section{Introduction}
Pokang\'a (Tiquié Pokang\'a) is described as a moribund Eastern Tucanoan variety of the Upper Tiqui\'e river region in northwest Brazil. Glottolog argues that the \texttt{pok} entry is spurious and should be merged with Waimaha (Northern Barasano, \texttt{bao}), since speakers say they speak Waimaha. Joshua Project lists about 200 Pokang\'a people in Brazil.

\section{Primary Educational Resources}
None found.

\section{Linguistic and Reference Materials}
\begin{itemize}
\item Glottolog entry (Tiqui\'e Pokang\'a): \url{https://glottolog.org/resource/languoid/iso/pok}
\item Wikipedia, Waimaj\~a language (the related Bar\'a/Waimaha language): \url{https://en.wikipedia.org/wiki/Waimaj%C3%A3_language}
\end{itemize}
No grammar or description specific to Pokang\'a was found.

\section{Religious and Community Resources}
\begin{itemize}
\item Joshua Project entry for Pokang\'a: \url{https://joshuaproject.net/index.php/languages/pok}
\end{itemize}

\section{Digital and Audio Resources}
None found.

\section{Curriculum Development Context}
No evidence of schooling in this variety was found. Any curriculum work should first clarify with the community whether materials in Waimaha/Bar\'a (Northern Barasano) are the appropriate basis.
\end{document}